% !TeX root = proyecto.tex
%=========================================================
\chapter{DISEÑO DEL SISTEMA}
\section{ARQUITECTURA GENERAL DEL SISTEMA}
\subsection{VISIÓN GENERAL}
El Sistema de Gestión de Cocina se diseñó bajo un modelo de arquitectura de tres capas, ejecutado mediante un despliegue local (on-premise) en las instalaciones físicas del establecimiento. Esta configuración se seleccionó con base en tres requerimientos críticos: la optimización de la latencia en la actualización de estados del Sistema de Visualización de Cocina, la necesidad de garantizar la operatividad continua ante interrupciones en la conectividad a la red de área amplia y el cumplimiento de las políticas de segmentación de red para el personal operativo. 

La infraestructura consolidó la lógica de negocio, el almacenamiento persistente de datos y los servicios de comunicación en tiempo real en un servidor físico central. El acceso a dichos recursos se gestionó mediante interfaces diferenciadas: una red de intranet dedicada para los módulos operativos (Gestión, Caja, Cocina y Entrega), y un puerto de acceso expuesto bajo configuraciones de seguridad específicas para el Módulo Cliente.
\begin{figure}[H]
    \centering
    \includesvg[width=0.4\textwidth, inkscapelatex=false]{diagrams/ArquitecturaGeneral/Arquitectura General.svg}
    \caption{Arquitectura del sistema completo.}
    \label{fig:arquitectura_sistema_completo}
\end{figure}

\subsection{STACK TENOLÓGICO}
\begingroup
\setlength{\LTleft}{\fill}
\setlength{\LTright}{\fill}

\begin{tabularx}{\textwidth}{l X X}
\caption{\textsc{Tabla de stack tecnológico.}} \label{tab:stack_tecnologico}\\
\toprule
\textbf{Capa} & \textbf{Tecnología} & \textbf{Justificación} \\
\midrule
\endfirsthead

\toprule
\textbf{Capa} & \textbf{Tecnología} & \textbf{Justificación} \\
\midrule
\endhead

\midrule
\multicolumn{3}{r}{\small Continúa en la página siguiente...} \\
\endfoot
\bottomrule
\endlastfoot

Frontend & Next.js + Tailwind CSS & SPA con componentes reutilizables, amplia comunidad y documentación, código abierto. \\
\addlinespace
Backend & Node.js + Express.js & Manejo eficiente de I / O asíncrono, ideal para WebSocket y API REST de alta concurrencia. \\
\addlinespace
Base de datos & PostgreSQL & RDBMS robusto, soporte nativo para JSON, transacciones ACID y código abierto. \\
\addlinespace
Tiempo real & Socket.io & Abstracción de WebSocket con fallback automático a long-polling. \\
\addlinespace
Autenticación & JWT + bcrypt & Estándar de autenticación sin estado, escalable y compatible con SPA. \\
\addlinespace
Contenerización & Docker + Docker Compose & Reproducibilidad del entorno, facilidad de despliegue y escalamiento. \\
\end{tabularx}
\endgroup

\subsection{HARDWARE}
Dado que es de interés que las cafeterías puedan ejecutar el software con equipo accesible y con el que puedan contar como equipos comerciales, se considerarán los siguientes requisitos mínimos para ejecutar el sistema:

\subsubsection{SERVIDOR CENTRAL}
\begingroup
\setlength{\LTleft}{\fill}
\setlength{\LTright}{\fill}

\begin{tabularx}{\textwidth}{l X X}
\caption{\textsc{Requisitos mínimos y recomendados para ejecutar el servidor central.}} \label{tab:requisitos_minimos_servidor_central}\\
\toprule
\textbf{Componente} & \textbf{Mínimo} & \textbf{Recomendado} \\
\midrule
\endfirsthead

\toprule
\textbf{Componente} & \textbf{Mínimo} & \textbf{Recomendado} \\
\midrule
\endhead

\midrule
\multicolumn{3}{r}{\small Continúa en la página siguiente...} \\
\endfoot
\bottomrule
\endlastfoot

CPU & Dual-core 2.0 GHz & Quad-core 2.5 GHz \\
\addlinespace
RAM & 4 GB & 8 GB \\
\addlinespace
Almacenamiento & 60 GB HDD & 120 GB SSD \\
\addlinespace
Red & Ethernet 100 Mbps & Ethernet + 10 Gbps \\
\addlinespace
OS & Ubuntu Server 22.04 LTS & Ubuntu Server 22.04 LTS \\
\end{tabularx}
\endgroup

Se recomienda almacenamiento SSD para reducir la latencia de escritura y lectura de la BD. Le red recomendada, por otro lado, permite la transferencia rápida de datos para la diversidad de estaciones conectadas a la vez. 

\subsubsection{ESTACIONES FIJAS}
\begingroup
\setlength{\LTleft}{\fill}
\setlength{\LTright}{\fill}

\begin{tabularx}{\textwidth}{l X}
\caption{\textsc{Requisitos mínimos para ejecución del sistema en estaciones fijas.}} \label{tab:requisitos_minimos_estaciones_fijas}\\
\toprule
\textbf{Componente} & \textbf{Mínimo} \\
\midrule
\endfirsthead

\toprule
\textbf{Componente} & \textbf{Mínimo} \\
\midrule
\endhead

\midrule
\multicolumn{2}{r}{\small Continúa en la página siguiente...} \\
\endfoot
\bottomrule
\endlastfoot

CPU & Dual-core 1.8 GHz \\
\addlinespace
RAM & 4 GB \\
\addlinespace
Almacenamiento & 32 GB \\
\addlinespace
Red & Ethernet \\
\addlinespace
OS & Ubuntu 22.04 \\
\end{tabularx}
\endgroup

Estos requisitos se encuentran cubiertos por la obsolescencia tecnológica actual por lo que las cafeterías contarán con alternativas accesibles para obtener el hardware necesario de operación en soluciones comerciales disponibles de las cuales poseerán en su totalidad. 

\section{DISEÑO DE LA BASE DE DATOS}
El diseño de la base de datos del \textit{Kitchen Management System} (KMS) se concibió como un esquema relacional centralizado, implementado sobre PostgreSQL, cuya responsabilidad es mantener la coherencia transaccional del ciclo completo de un pedido: desde su captura en caja hasta su entrega al cliente. La decisión de adoptar un motor relacional responde al requisito RFUN-05, que establece una base de datos única y centralizada para todos los módulos del sistema, garantizando integridad referencial y eliminando inconsistencias derivadas de almacenes de datos distribuidos.

El modelo se organizó en nueve dominios funcionales: (1)~autenticación y personal, (2)~catálogo de productos e insumos, (3)~recetas y pasos de preparación, (4)~clientes y ciclo de vida del pedido, (5)~instantáneas de producción para el KDS (\textit{Kitchen Display System}), (6)~entrega y cancelación, (7)~cortes de caja e inventario, (8)~configuración del sistema y (9)~auditoría. Esta segmentación favoreció la aplicación sistemática de las reglas de normalización y facilitó el análisis de dependencias funcionales entre atributos.

Un principio de diseño transversal fue la \textit{inmutabilidad por instantánea}: cuando un pedido se confirma, los datos operativos de la receta (pasos, ingredientes, duraciones, estación asignada) se copian a tablas de producción independientes (\texttt{PASO\_RECETA\_PEDIDO}, \texttt{INGREDIENTE\_PASO\_PEDIDO}). Esto garantiza que modificaciones posteriores al catálogo no alteren pedidos en curso ni el historial de producción, satisfaciendo las reglas de negocio RN-COC-001 y RN-CAJ-007.

\subsection{DICCIONARIO DE DATOS}
\subsubsection{USUARIO}
\begingroup
\setlength{\LTleft}{\fill}
\setlength{\LTright}{\fill}

\begin{tabularx}{\textwidth}{p{3cm} p{0.15\textwidth} p{0.15\textwidth} X}
\caption{\textsc{Tabla USUARIO de la base de datos.}} \label{tab:diccionario-usuario}\\
\toprule
\textbf{Nombre del campo} & \textbf{Tipo de dato} & \textbf{Restricciones} & \textbf{Descripción funcional} \\
\midrule
\endfirsthead

\toprule
\textbf{Nombre del campo} & \textbf{Tipo de dato} & \textbf{Restricciones} & \textbf{Descripción funcional} \\
\midrule
\endhead

\midrule
\multicolumn{4}{r}{\small Continúa en la página siguiente...} \\
\endfoot
\bottomrule
\endlastfoot

\texttt{id\_usuario} & UUID & PK, NOT NULL & Identificador único del usuario con acceso al panel de Manager o Caja (UUID v4). \\
\addlinespace
\texttt{nombre\_usuario} & VARCHAR(50) & UNIQUE, NOT NULL & Credencial de inicio de sesión. Debe ser única en el sistema. \\
\addlinespace
\texttt{nombre\_completo} & VARCHAR(150) & NOT NULL & Nombre real del usuario para visualización en bitácoras y reportes. \\
\addlinespace
\texttt{hash\_contrasena} & TEXT & NOT NULL & Hash bcrypt de la contraseña. Nunca se almacena en texto plano. \\
\addlinespace
\texttt{rol} & VARCHAR(20) & NOT NULL, CHECK (MANAGER\textbar CAJERO) & Rol que determina los privilegios RBAC del usuario dentro del sistema (RNFU-01). \\
\addlinespace
\texttt{activo} & BOOLEAN & NOT NULL, DEFAULT TRUE & Bandera de estado lógico; los usuarios se desactivan en lugar de eliminarse para preservar integridad referencial. \\
\addlinespace
\texttt{fecha\_creacion} & TIMESTAMP WITH TIME ZONE & NOT NULL, DEFAULT NOW() & Marca temporal del momento de creación del registro. \\
\end{tabularx}
\endgroup

\subsubsection{EMPLEADO}
\begingroup
\setlength{\LTleft}{\fill}
\setlength{\LTright}{\fill}

\begin{tabularx}{\textwidth}{p{3cm} p{2.5cm} p{3cm} X}
\caption{\textsc{Tabla EMPLEADO de la base de datos.}} \label{tab:diccionario-empleado}\\
\toprule
\textbf{Nombre del campo} & \textbf{Tipo de dato} & \textbf{Restricciones} & \textbf{Descripción funcional} \\
\midrule
\endfirsthead

\toprule
\textbf{Nombre del campo} & \textbf{Tipo de dato} & \textbf{Restricciones} & \textbf{Descripción funcional} \\
\midrule
\endhead

\midrule
\multicolumn{4}{r}{\small Continúa en la página siguiente...} \\
\endfoot
\bottomrule
\endlastfoot

\texttt{id\_empleado} & UUID & PK, NOT NULL & Identificador único del empleado de cocina (UUID v4). \\
\addlinespace
\texttt{id\_usuario} & UUID & FK $\rightarrow$ USUARIO, UNIQUE, NULL & Referencia al usuario del sistema si el empleado también tiene acceso al panel Manager/Caja. Nullable para cocineros sin acceso a panel. \\
\addlinespace
\texttt{nombre\_completo} & VARCHAR(150) & NOT NULL & Nombre del cocinero para asignación de turnos y visualización en el KDS. \\
\addlinespace
\texttt{nip\_hash} & TEXT & NOT NULL & Hash bcrypt del NIP de 4-6 dígitos usado para autenticación rápida en el KDS. \\
\addlinespace
\texttt{activo} & BOOLEAN & NOT NULL, DEFAULT TRUE & Estado lógico del empleado. Los empleados inactivos no se pueden asignar a turnos. \\
\end{tabularx}
\endgroup

\subsubsection{TURNO}
\begingroup
\setlength{\LTleft}{\fill}
\setlength{\LTright}{\fill}

\begin{tabularx}{\textwidth}{l p{2.5cm} p{3cm} X}
\caption{\textsc{Tabla TURNO de la base de datos.}} \label{tab:turno}\\
\toprule
\textbf{Nombre del campo} & \textbf{Tipo de dato} & \textbf{Restricciones} & \textbf{Descripción funcional} \\
\midrule
\endfirsthead

\toprule
\textbf{Nombre del campo} & \textbf{Tipo de dato} & \textbf{Restricciones} & \textbf{Descripción funcional} \\
\midrule
\endhead

\midrule
\multicolumn{4}{r}{\small Continúa en la página siguiente...} \\
\endfoot
\bottomrule
\endlastfoot

\texttt{id\_turno} & UUID & PK, NOT NULL & Identificador único del turno de trabajo (UUID v4). \\
\addlinespace
\texttt{hora\_inicio} & TIME & NOT NULL & Hora de inicio del turno. Validada contra \texttt{hora\_fin} en regla RN-MGR-026. \\
\addlinespace
\texttt{hora\_fin} & TIME & NOT NULL & Hora de fin del turno. Debe ser posterior a \texttt{hora\_inicio}. \\
\addlinespace
\texttt{descripcion} & VARCHAR(100) & NOT NULL & Nombre descriptivo del turno (ej. ``Turno Mañana'', ``Turno Tarde''). \\
\end{tabularx}
\endgroup

\subsubsection{ESTACION}
\begingroup
\setlength{\LTleft}{\fill}
\setlength{\LTright}{\fill}

\begin{tabularx}{\textwidth}{l p{2.5cm} p{3cm} X}
\caption{\textsc{Tabla ESTACION de la base de datos.}} \label{tab:estacion}\\
\toprule
\textbf{Nombre del campo} & \textbf{Tipo de dato} & \textbf{Restricciones} & \textbf{Descripción funcional} \\
\midrule
\endfirsthead

\toprule
\textbf{Nombre del campo} & \textbf{Tipo de dato} & \textbf{Restricciones} & \textbf{Descripción funcional} \\
\midrule
\endhead

\midrule
\multicolumn{4}{r}{\small Continúa en la página siguiente...} \\
\endfoot
\bottomrule
\endlastfoot

\texttt{id\_estacion} & UUID & PK, NOT NULL & Identificador único de la estación de trabajo (UUID v4). \\
\addlinespace
\texttt{nombre} & VARCHAR(100) & NOT NULL & Nombre operativo de la estación (ej. ``Caja Principal'', ``Parrilla'', ``Módulo Entrega''). \\
\addlinespace
\texttt{tipo} & VARCHAR(20) & NOT NULL, CHECK (CAJA  $|$  COCINA  $|$  ENTREGA) & Determina el módulo de interfaz que se sirve a ese dispositivo y los privilegios de acceso. \\
\addlinespace
\texttt{activa} & BOOLEAN & NOT NULL, DEFAULT TRUE & Estado lógico. Las estaciones inactivas no reciben asignaciones ni pedidos. \\
\end{tabularx}
\endgroup

\subsubsection{CREDENCIAL\_ESTACION}

\begingroup
\setlength{\LTleft}{\fill}
\setlength{\LTright}{\fill}

\begin{tabularx}{\textwidth}{p{3cm} p{2.5cm} p{3.5cm} X}
\caption{\textsc{Tabla CREDENCIAL\_ESTACION de la base de datos.}} \label{tab:credencial-estacion}\\
\toprule
\textbf{Nombre del campo} & \textbf{Tipo de dato} & \textbf{Restricciones} & \textbf{Descripción funcional} \\
\midrule
\endfirsthead

\toprule
\textbf{Nombre del campo} & \textbf{Tipo de dato} & \textbf{Restricciones} & \textbf{Descripción funcional} \\
\midrule
\endhead

\midrule
\multicolumn{4}{r}{\small Continúa en la página siguiente...} \\
\endfoot
\bottomrule
\endlastfoot

\texttt{id\_credencial} & UUID & PK, NOT NULL & Identificador único de la credencial compartida (UUID v4). \\
\addlinespace
\texttt{id\_estacion} & UUID & FK $\rightarrow$ ESTACION, NOT NULL & Estación a la que pertenece esta credencial. \\
\addlinespace
\texttt{nombre\_usuario} & VARCHAR(50) & UNIQUE, NOT NULL & Usuario de inicio de sesión compartido para el puesto físico (Caja o Entrega). \\
\addlinespace
\texttt{hash\_contrasena} & TEXT & NOT NULL & Hash bcrypt de la contraseña compartida del puesto. \\
\addlinespace
\texttt{activa} & BOOLEAN & NOT NULL, DEFAULT TRUE & Estado lógico de la credencial. \\
\end{tabularx}
\endgroup

\subsubsection{ASIGNACION\_ESTACION}
\begingroup
\setlength{\LTleft}{\fill}
\setlength{\LTright}{\fill}

\begin{tabularx}{\textwidth}{p{3cm} p{2.5cm} p{3.5cm} X}
\caption{\textsc{Tabla ASIGNACION\_ESTACION de la base de datos.}} \label{tab:asignacion-estacion}\\
\toprule
\textbf{Nombre del campo} & \textbf{Tipo de dato} & \textbf{Restricciones} & \textbf{Descripción funcional} \\
\midrule
\endfirsthead

\toprule
\textbf{Nombre del campo} & \textbf{Tipo de dato} & \textbf{Restricciones} & \textbf{Descripción funcional} \\
\midrule
\endhead

\midrule
\multicolumn{4}{r}{\small Continúa en la página siguiente...} \\
\endfoot
\bottomrule
\endlastfoot

\texttt{id\_asignacion} & UUID & PK, NOT NULL & Identificador único de la asignación (UUID v4). \\
\addlinespace
\texttt{id\_empleado} & UUID & FK $\rightarrow$ EMPLEADO, NOT NULL & Empleado asignado. \\
\addlinespace
\texttt{id\_estacion} & UUID & FK $\rightarrow$ ESTACION, NOT NULL & Estación a la que se asigna el empleado. Debe ser de tipo COCINA. \\
\addlinespace
\texttt{id\_turno} & UUID & FK $\rightarrow$ TURNO, NOT NULL & Turno que cubre esta asignación. \\
\addlinespace
\texttt{fecha} & DATE & NOT NULL & Fecha calendario de la asignación. \\
\addlinespace
\texttt{inicio\_sesion} & TIMESTAMP WITH TIME ZONE & NULL & Marca temporal del momento en que el empleado autenticó en el KDS. \\
\addlinespace
\texttt{fin\_sesion} & TIMESTAMP WITH TIME ZONE & NULL & Marca temporal del cierre de sesión. Nullable mientras la sesión está activa. \\
\end{tabularx}
\endgroup

\subsubsection{INSUMO}
\begingroup
\setlength{\LTleft}{\fill}
\setlength{\LTright}{\fill}

\begin{tabularx}{\textwidth}{p{3cm} p{2.5cm} p{3.5cm} X}
\caption{\textsc{Tabla INSUMO de la base de datos.}} \label{tab:insumo}\\
\toprule
\textbf{Nombre del campo} & \textbf{Tipo de dato} & \textbf{Restricciones} & \textbf{Descripción funcional} \\
\midrule
\endfirsthead

\toprule
\textbf{Nombre del campo} & \textbf{Tipo de dato} & \textbf{Restricciones} & \textbf{Descripción funcional} \\
\midrule
\endhead

\midrule
\multicolumn{4}{r}{\small Continúa en la página siguiente...} \\
\endfoot
\bottomrule
\endlastfoot

\texttt{id\_insumo} & UUID & PK, NOT NULL & Identificador único del insumo (UUID v4). \\
\addlinespace
\texttt{nombre} & VARCHAR(150) & NOT NULL & Denominación del insumo (ej. ``Harina de trigo'', ``Vaso desechable 16 oz''). \\
\addlinespace
\texttt{unidad\_medida} & VARCHAR(30) & NOT NULL & Unidad en que se expresa el stock (ej. ``kg'', ``L'', ``piezas''). \\
\addlinespace
\texttt{stock\_actual} & DECIMAL(10,3) & NOT NULL, DEFAULT 0 & Cantidad disponible en inventario en el momento de la consulta. \\
\addlinespace
\texttt{stock\_minimo} & DECIMAL(10,3) & NOT NULL, DEFAULT 0 & Umbral mínimo por debajo del cual se genera una alerta de reabastecimiento. \\
\addlinespace
\texttt{activo} & BOOLEAN & NOT NULL, DEFAULT TRUE & Estado lógico del insumo. Los insumos inactivos no se asignan a nuevas recetas. \\
\end{tabularx}
\endgroup

\subsubsection{PRODUCTO}
\begingroup
\setlength{\LTleft}{\fill}
\setlength{\LTright}{\fill}

\begin{tabularx}{\textwidth}{p{4cm} p{2.5cm} p{3cm} X}
\caption{\textsc{Tabla PRODUCTO de la base de datos.}} \label{tab:producto}\\
\toprule
\textbf{Nombre del campo} & \textbf{Tipo de dato} & \textbf{Restricciones} & \textbf{Descripción funcional} \\
\midrule
\endfirsthead

\toprule
\textbf{Nombre del campo} & \textbf{Tipo de dato} & \textbf{Restricciones} & \textbf{Descripción funcional} \\
\midrule
\endhead

\midrule
\multicolumn{4}{r}{\small Continúa en la página siguiente...} \\
\endfoot
\bottomrule
\endlastfoot

\texttt{id\_producto} & UUID & PK, NOT NULL & Identificador único del producto del menú (UUID v4). \\
\addlinespace
\texttt{nombre} & VARCHAR(150) & NOT NULL & Nombre del artículo tal como aparece en el menú del cliente. \\
\addlinespace
\texttt{descripcion} & TEXT & NULL & Descripción breve para el módulo de cliente. \\
\addlinespace
\texttt{precio\_base} & DECIMAL(10,2) & NOT NULL, CHECK ($\ge$ 0) & Precio de venta estándar antes de aplicar promociones. \\
\addlinespace
\texttt{disponible} & BOOLEAN & NOT NULL, DEFAULT TRUE & Disponibilidad operativa. Un producto no disponible no aparece en el menú del cliente. \\
\addlinespace
\texttt{permite\_personalizacion} & BOOLEAN & NOT NULL, DEFAULT FALSE & Indica si el cliente puede modificar ingredientes al agregar el producto al carrito. \\
\end{tabularx}
\endgroup

\subsubsection{PASO\_RECETA}
\begingroup
\setlength{\LTleft}{\fill}
\setlength{\LTright}{\fill}

\begin{tabularx}{\textwidth}{p{4cm} p{2.5cm} p{3cm} X}
\caption{\textsc{Tabla PASO\_RECETA de la base de datos.}} \label{tab:paso-receta}\\
\toprule
\textbf{Nombre del campo} & \textbf{Tipo de dato} & \textbf{Restricciones} & \textbf{Descripción funcional} \\
\midrule
\endfirsthead

\toprule
\textbf{Nombre del campo} & \textbf{Tipo de dato} & \textbf{Restricciones} & \textbf{Descripción funcional} \\
\midrule
\endhead

\midrule
\multicolumn{4}{r}{\small Continúa en la página siguiente...} \\
\endfoot
\bottomrule
\endlastfoot

\texttt{id\_paso\_receta} & UUID & PK, NOT NULL & Identificador único del paso de receta (UUID v4). \\
\addlinespace
\texttt{id\_producto} & UUID & FK $\rightarrow$ PRODUCTO, NOT NULL & Producto al que pertenece este paso. \\
\addlinespace
\texttt{id\_estacion} & UUID & FK $\rightarrow$ ESTACION, NOT NULL & Estación de cocina responsable de ejecutar el paso. Determina en qué KDS aparece. \\
\addlinespace
\texttt{nombre} & VARCHAR(150) & NOT NULL & Denominación del paso (ej. ``Asar la carne'', ``Ensamblar el sándwich''). \\
\addlinespace
\texttt{orden} & INTEGER & NOT NULL, CHECK ($\ge$ 1) & Posición secuencial sugerida dentro de la receta, complementada por las dependencias en DAG. \\
\addlinespace
\texttt{duracion\_estimada\_seg} & INTEGER & NOT NULL, CHECK ($>$ 0) & Duración esperada del paso en segundos. Alimenta el cálculo de urgencia en el KDS. \\
\addlinespace
\texttt{activo} & BOOLEAN & NOT NULL, DEFAULT TRUE & Estado lógico. Los pasos inactivos no se copian a nuevas instantáneas de producción. \\
\end{tabularx}
\endgroup

\subsubsection{CLIENTE}
\begingroup
\setlength{\LTleft}{\fill}
\setlength{\LTright}{\fill}

\begin{tabularx}{\textwidth}{p{3cm} p{2.5cm} p{3.5cm} X}
\caption{\textsc{Tabla CLIENTE de la base de datos.}} \label{tab:cliente}\\
\toprule
\textbf{Nombre del campo} & \textbf{Tipo de dato} & \textbf{Restricciones} & \textbf{Descripción funcional} \\
\midrule
\endfirsthead

\toprule
\textbf{Nombre del campo} & \textbf{Tipo de dato} & \textbf{Restricciones} & \textbf{Descripción funcional} \\
\midrule
\endhead

\midrule
\multicolumn{4}{r}{\small Continúa en la página siguiente...} \\
\endfoot
\bottomrule
\endlastfoot

\texttt{id\_cliente} & UUID & PK, NOT NULL & Identificador único del cliente registrado (UUID v4). \\
\addlinespace
\texttt{correo} & VARCHAR(254) & UNIQUE, NOT NULL & Dirección de correo electrónico usada como identificador de cuenta y para notificaciones SMTP. \\
\addlinespace
\texttt{nombre} & VARCHAR(150) & NOT NULL & Nombre del cliente para personalización de la interfaz. \\
\addlinespace
\texttt{hash\_contrasena} & TEXT & NOT NULL & Hash bcrypt de la contraseña del cliente. \\
\addlinespace
\texttt{cuenta\_verificada} & BOOLEAN & NOT NULL, DEFAULT FALSE & Estado de verificación. La cuenta debe verificarse por correo antes de poder realizar pedidos. \\
\addlinespace
\texttt{fecha\_registro} & TIMESTAMP WITH TIME ZONE & NOT NULL, DEFAULT NOW() & Marca temporal del registro. \\
\end{tabularx}
\endgroup

\subsubsection{PEDIDO}
\begingroup
\setlength{\LTleft}{\fill}
\setlength{\LTright}{\fill}

\begin{tabularx}{\textwidth}{p{3cm} p{2.5cm} p{3.5cm} X}
\caption{\textsc{Tabla PEDIDO de la base de datos.}} \label{tab:pedido}\\
\toprule
\textbf{Nombre del campo} & \textbf{Tipo de dato} & \textbf{Restricciones} & \textbf{Descripción funcional} \\
\midrule
\endfirsthead

\toprule
\textbf{Nombre del campo} & \textbf{Tipo de dato} & \textbf{Restricciones} & \textbf{Descripción funcional} \\
\midrule
\endhead

\midrule
\multicolumn{4}{r}{\small Continúa en la página siguiente...} \\
\endfoot
\bottomrule
\endlastfoot

\texttt{id\_pedido} & UUID & PK, NOT NULL & Identificador único del pedido (UUID v4). Se usa como base para el token QR. \\
\addlinespace
\texttt{id\_cliente} & UUID & FK $\rightarrow$ CLIENTE, NULL & Cliente que realizó el pedido. NULL si el pedido es anónimo (\texttt{es\_anonimo} = TRUE). \\
\addlinespace
\texttt{id\_estacion} & UUID & FK $\rightarrow$ ESTACION, NOT NULL & Estación de caja en la que se confirmó el pedido. \\
\addlinespace
\texttt{estado} & VARCHAR(20) & NOT NULL, CHECK (EN\_PREPARACION  $|$  PREPARADO  $|$  ENTREGADO  $|$  CANCELADO) & Estado actual en la máquina de estados del pedido. \\
\addlinespace
\texttt{total} & DECIMAL(10,2) & NOT NULL, CHECK ($\ge$ 0) & Monto total cobrado al cliente, calculado como suma de \texttt{ITEM\_PEDIDO.subtotal}. \\
\addlinespace
\texttt{metodo\_pago} & VARCHAR(50) & NOT NULL & Método de pago registrado al confirmar (ej. ``EFECTIVO'', ``TARJETA''). \\
\addlinespace
\texttt{es\_anonimo} & BOOLEAN & NOT NULL, DEFAULT FALSE & Indica si el pedido fue realizado sin cuenta de cliente. \\
\addlinespace
\texttt{ventana\_inicio} & TIME & NULL & Inicio de la ventana de recolección programada. Solo aplica cuando el sistema tiene ventanas activas. \\
\addlinespace
\texttt{ventana\_fin} & TIME & NULL & Fin de la ventana de recolección programada. \\
\addlinespace
\texttt{fecha\_creacion} & TIMESTAMP WITH TIME ZONE & NOT NULL, DEFAULT NOW() & Marca temporal de creación del pedido. \\
\end{tabularx}
\endgroup

\subsubsection{ITEM\_PEDIDO}
\begingroup
\setlength{\LTleft}{\fill}
\setlength{\LTright}{\fill}

\begin{tabularx}{\textwidth}{p{4cm} p{2.5cm} p{3cm} X}
\caption{\textsc{Tabla ITEM\_PEDIDO de la base de datos.}} \label{tab:item-pedido}\\
\toprule
\textbf{Nombre del campo} & \textbf{Tipo de dato} & \textbf{Restricciones} & \textbf{Descripción funcional} \\
\midrule
\endfirsthead

\toprule
\textbf{Nombre del campo} & \textbf{Tipo de dato} & \textbf{Restricciones} & \textbf{Descripción funcional} \\
\midrule
\endhead

\midrule
\multicolumn{4}{r}{\small Continúa en la página siguiente...} \\
\endfoot
\bottomrule
\endlastfoot

\texttt{id\_item\_pedido} & UUID & PK, NOT NULL & Identificador único del ítem dentro del pedido (UUID v4). \\
\addlinespace
\texttt{id\_pedido} & UUID & FK $\rightarrow$ PEDIDO, NOT NULL & Pedido al que pertenece el ítem. \\
\addlinespace
\texttt{id\_producto} & UUID & FK $\rightarrow$ PRODUCTO, NOT NULL & Referencia trazable al catálogo de productos. \\
\addlinespace
\texttt{cantidad} & INTEGER & NOT NULL, CHECK ($\ge$ 1) & Número de unidades del producto incluidas en este ítem. \\
\addlinespace
\texttt{precio\_unitario\_snapshot} & DECIMAL(10,2) & NOT NULL & Precio del producto en el momento de confirmación del pedido. Inmutable por diseño (snapshot). \\
\addlinespace
\texttt{subtotal} & DECIMAL(10,2) & NOT NULL & Resultado de $\text{cantidad} \times \text{precio\_unitario\_snapshot}$. Calculado y persistido para eficiencia en reportes. \\
\addlinespace
\texttt{notas} & TEXT & NULL & Instrucciones especiales del cliente no capturables con los ingredientes parametrizados. \\
\end{tabularx}
\endgroup

\subsubsection{INGREDIENTE\_ITEM\_PEDIDO}
\begingroup
\setlength{\LTleft}{\fill}
\setlength{\LTright}{\fill}

\begin{tabularx}{\textwidth}{p{3cm} p{2.5cm} p{3.5cm} X}
\caption{\textsc{Tabla INGREDIENTE\_ITEM\_PEDIDO de la base de datos.}} \label{tab:ingrediente-item-pedido}\\
\toprule
\textbf{Nombre del campo} & \textbf{Tipo de dato} & \textbf{Restricciones} & \textbf{Descripción funcional} \\
\midrule
\endfirsthead

\toprule
\textbf{Nombre del campo} & \textbf{Tipo de dato} & \textbf{Restricciones} & \textbf{Descripción funcional} \\
\midrule
\endhead

\midrule
\multicolumn{4}{r}{\small Continúa en la página siguiente...} \\
\endfoot
\bottomrule
\endlastfoot

\texttt{id\_ingrediente\_item} & UUID & PK, NOT NULL & Identificador único de la personalización del ítem (UUID v4). \\
\addlinespace
\texttt{id\_item\_pedido} & UUID & FK $\rightarrow$ ITEM\_PEDIDO, NOT NULL & Ítem al que aplica esta personalización. \\
\addlinespace
\texttt{id\_insumo} & UUID & FK $\rightarrow$ INSUMO, NOT NULL & Insumo personalizado (solo los marcados como \texttt{es\_personalizable} en la receta). \\
\addlinespace
\texttt{cantidad} & DECIMAL(10,3) & NOT NULL, CHECK ($\ge$ 0) & Cantidad del insumo a usar. Cero cuando la acción es EXCLUIR. \\
\addlinespace
\texttt{accion} & VARCHAR(10) & NOT NULL, CHECK (INCLUIR $|$ EXCLUIR $|$ EXTRA) & Tipo de modificación solicitada por el cliente. \\
\end{tabularx}
\endgroup

\subsubsection{CODIGO\_VERIFICACION\_QR}
\begingroup
\setlength{\LTleft}{\fill}
\setlength{\LTright}{\fill}

\begin{tabularx}{\textwidth}{p{3cm} p{2.5cm} p{3.5cm} X}
\caption{\textsc{Tabla CODIGO\_VERIFICACION\_QR de la base de datos.}} \label{tab:codigo-verificacion-qr}\\
\toprule
\textbf{Nombre del campo} & \textbf{Tipo de dato} & \textbf{Restricciones} & \textbf{Descripción funcional} \\
\midrule
\endfirsthead

\toprule
\textbf{Nombre del campo} & \textbf{Tipo de dato} & \textbf{Restricciones} & \textbf{Descripción funcional} \\
\midrule
\endhead

\midrule
\multicolumn{4}{r}{\small Continúa en la página siguiente...} \\
\endfoot
\bottomrule
\endlastfoot

\texttt{id\_codigo\_qr} & UUID & PK, NOT NULL & Identificador único del código QR (UUID v4). \\
\addlinespace
\texttt{id\_pedido} & UUID & FK $\rightarrow$ PEDIDO, UNIQUE, NOT NULL & Pedido al que corresponde el QR. Relación 1:1; un pedido tiene un único QR vigente. \\
\addlinespace
\texttt{token} & VARCHAR(512) & UNIQUE, NOT NULL & Token opaco generado según RNFU-06 (UUID v4 o JWT [INFERENCIA: formato exacto no especificado]). \\
\addlinespace
\texttt{estado} & VARCHAR(10) & NOT NULL, CHECK (ACTIVO  $|$  CONSUMIDO  $|$  EXPIRADO) & Ciclo de vida del token. Solo los tokens ACTIVOS son válidos para entrega. \\
\addlinespace
\texttt{fecha\_generacion} & TIMESTAMP WITH TIME ZONE & NOT NULL & Momento en que se generó el token. \\
\addlinespace
\texttt{fecha\_expiracion} & TIMESTAMP WITH TIME ZONE & NOT NULL & Tiempo límite de validez del token (calculado al generar). \\
\addlinespace
\texttt{fecha\_uso} & TIMESTAMP WITH TIME ZONE & NULL & Momento en que el token fue consumido por el módulo de Entrega. NULL si aún no se usa. \\
\end{tabularx}
\endgroup

\subsubsection{PEDIDO\_ACCESO\_ANONIMO}

\begingroup
\setlength{\LTleft}{\fill}
\setlength{\LTright}{\fill}

\begin{tabularx}{\textwidth}{l p{2.5cm} p{3.5cm} X}
\caption{\textsc{Tabla PEDIDO\_ACCESO\_ANONIMO de la base de datos.}} \label{tab:pedido-acceso-anonimo}\\
\toprule
\textbf{Nombre del campo} & \textbf{Tipo de dato} & \textbf{Restricciones} & \textbf{Descripción funcional} \\
\midrule
\endfirsthead

\toprule
\textbf{Nombre del campo} & \textbf{Tipo de dato} & \textbf{Restricciones} & \textbf{Descripción funcional} \\
\midrule
\endhead

\midrule
\multicolumn{4}{r}{\small Continúa en la página siguiente...} \\
\endfoot
\bottomrule
\endlastfoot

\texttt{id\_acceso\_anonimo} & UUID & PK, NOT NULL & Identificador único del acceso anónimo (UUID v4). \\
\addlinespace
\texttt{id\_pedido} & UUID & FK $\rightarrow$ PEDIDO, UNIQUE, NOT NULL & Pedido cuyo seguimiento se expone de forma anónima. \\
\addlinespace
\texttt{token} & VARCHAR(512) & UNIQUE, NOT NULL & Token opaco de acceso de un solo uso, entregado al cliente anónimo en caja. \\
\addlinespace
\texttt{fecha\_creacion} & TIMESTAMP WITH TIME ZONE & NOT NULL & Momento de generación del token de acceso anónimo. \\
\addlinespace
\texttt{fecha\_expiracion} & TIMESTAMP WITH TIME ZONE & NOT NULL & Tiempo límite de validez. El token expira cuando el pedido se entrega o cancela. \\
\end{tabularx}
\endgroup

\subsubsection{PEDIDO\_ESTADO\_TIMESTAMP}
\begingroup
\setlength{\LTleft}{\fill}
\setlength{\LTright}{\fill}

\begin{tabularx}{\textwidth}{l p{2.5cm} p{3.5cm} X}
\caption{\textsc{Tabla PEDIDO\_ESTADO\_TIMESTAMP de la base de datos.}} \label{tab:pedido-estado-timestamp}\\
\toprule
\textbf{Nombre del campo} & \textbf{Tipo de dato} & \textbf{Restricciones} & \textbf{Descripción funcional} \\
\midrule
\endfirsthead

\toprule
\textbf{Nombre del campo} & \textbf{Tipo de dato} & \textbf{Restricciones} & \textbf{Descripción funcional} \\
\midrule
\endhead

\midrule
\multicolumn{4}{r}{\small Continúa en la página siguiente...} \\
\endfoot
\bottomrule
\endlastfoot

\texttt{id\_timestamp} & UUID & PK, NOT NULL & Identificador único del registro de cambio de estado (UUID v4). \\
\addlinespace
\texttt{id\_pedido} & UUID & FK $\rightarrow$ PEDIDO, NOT NULL & Pedido cuyo estado cambió. \\
\addlinespace
\texttt{estado} & VARCHAR(20) & NOT NULL & Estado al que transitó el pedido en este momento. \\
\addlinespace
\texttt{id\_empleado} & UUID & FK $\rightarrow$ EMPLEADO, NULL & Empleado que realizó el cambio de estado. NULL si el cambio fue automático o iniciado por sistema. \\
\addlinespace
\texttt{fecha\_cambio} & TIMESTAMP WITH TIME ZONE & NOT NULL, DEFAULT NOW() & Marca temporal exacta del cambio de estado. Permite reconstruir el historial completo. \\
\end{tabularx}
\endgroup

\subsubsection{PASO\_RECETA\_PEDIDO}
\begingroup
\setlength{\LTleft}{\fill}
\setlength{\LTright}{\fill}

\begin{tabularx}{\textwidth}{l p{2.5cm} p{3.5cm} X}
\caption{\textsc{Tabla PASO\_RECETA\_PEDIDO de la base de datos.}} \label{tab:paso-receta-pedido}\\
\toprule
\textbf{Nombre del campo} & \textbf{Tipo de dato} & \textbf{Restricciones} & \textbf{Descripción funcional} \\
\midrule
\endfirsthead

\toprule
\textbf{Nombre del campo} & \textbf{Tipo de dato} & \textbf{Restricciones} & \textbf{Descripción funcional} \\
\midrule
\endhead

\midrule
\multicolumn{4}{r}{\small Continúa en la página siguiente...} \\
\endfoot
\bottomrule
\endlastfoot

\texttt{id\_paso\_receta\_pedido} & UUID & PK, NOT NULL & Identificador único del paso en la instantánea de producción (UUID v4). \\
\addlinespace
\texttt{id\_item\_pedido} & UUID & FK $\rightarrow$ ITEM\_PEDIDO, NOT NULL & Ítem del pedido al que pertenece este paso. \\
\addlinespace
\texttt{id\_paso\_receta} & UUID & FK $\rightarrow$ PASO\_RECETA, NOT NULL & Referencia trazable al paso del catálogo del que se generó la instantánea. \\
\addlinespace
\texttt{id\_estacion} & UUID & FK $\rightarrow$ ESTACION, NOT NULL & Estación de cocina asignada (copiada del catálogo en el momento de confirmación). \\
\addlinespace
\texttt{nombre\_snapshot} & VARCHAR(150) & NOT NULL & Nombre del paso en el momento de la confirmación del pedido. Inmutable. \\
\addlinespace
\texttt{duracion\_estimada\_seg\_snapshot} & INTEGER & NOT NULL & Duración estimada copiada del catálogo al confirmar. Inmutable. \\
\addlinespace
\texttt{orden\_snapshot} & INTEGER & NOT NULL & Orden sugerido copiado del catálogo al confirmar. \\
\end{tabularx}
\endgroup

\subsubsection{PASO\_PEDIDO\_ESTADO}
\begingroup
\setlength{\LTleft}{\fill}
\setlength{\LTright}{\fill}

\begin{tabularx}{\textwidth}{l p{2.5cm} p{3.5cm} X}
\caption{\textsc{Tabla PASO\_PEDIDO\_ESTADO de la base de datos.}} \label{tab:paso-pedido-estado}\\
\toprule
\textbf{Nombre del campo} & \textbf{Tipo de dato} & \textbf{Restricciones} & \textbf{Descripción funcional} \\
\midrule
\endfirsthead

\toprule
\textbf{Nombre del campo} & \textbf{Tipo de dato} & \textbf{Restricciones} & \textbf{Descripción funcional} \\
\midrule
\endhead

\midrule
\multicolumn{4}{r}{\small Continúa en la página siguiente...} \\
\endfoot
\bottomrule
\endlastfoot

\texttt{id\_paso\_pedido\_estado} & UUID & PK, NOT NULL & Identificador único del estado de ejecución del paso (UUID v4). \\
\addlinespace
\texttt{id\_paso\_receta\_pedido} & UUID & FK $\rightarrow$ PASO\_RECETA\_PEDIDO, UNIQUE, NOT NULL & Paso de producción al que corresponde este estado. Relación 1:1. \\
\addlinespace
\texttt{id\_empleado\_asignado} & UUID & FK $\rightarrow$ EMPLEADO, NULL & Cocinero que tomó el paso. NULL si el paso aún no fue tomado. \\
\addlinespace
\texttt{estado} & VARCHAR(15) & NOT NULL, CHECK (PENDIENTE $|$ EN\_PROCESO  $|$  COMPLETADO), DEFAULT PENDIENTE & Estado de ejecución del paso. Alimenta la lógica de liberación de dependencias DAG. \\
\addlinespace
\texttt{fecha\_inicio} & TIMESTAMP WITH TIME ZONE & NULL & Momento en que el cocinero marcó el paso como iniciado. \\
\addlinespace
\texttt{fecha\_fin} & TIMESTAMP WITH TIME ZONE & NULL & Momento en que el cocinero marcó el paso como completado. \\
\end{tabularx}
\endgroup

\subsubsection{CANCELACION\_ENTREGA}
\begingroup
\setlength{\LTleft}{\fill}
\setlength{\LTright}{\fill}

\begin{tabularx}{\textwidth}{l p{2.5cm} p{3.5cm} X}
\caption{\textsc{Tabla CANCELACION\_ENTREGA de la base de datos.}} \label{tab:cancelacion-entrega}\\
\toprule
\textbf{Nombre del campo} & \textbf{Tipo de dato} & \textbf{Restricciones} & \textbf{Descripción funcional} \\
\midrule
\endfirsthead

\toprule
\textbf{Nombre del campo} & \textbf{Tipo de dato} & \textbf{Restricciones} & \textbf{Descripción funcional} \\
\midrule
\endhead

\midrule
\multicolumn{4}{r}{\small Continúa en la página siguiente...} \\
\endfoot
\bottomrule
\endlastfoot

\texttt{id\_cancelacion} & UUID & PK, NOT NULL & Identificador único del registro de cancelación (UUID v4). \\
\addlinespace
\texttt{id\_pedido} & UUID & FK $\rightarrow$ PEDIDO, UNIQUE, NOT NULL & Pedido que fue cancelado. Un pedido solo puede cancelarse una vez. \\
\addlinespace
\texttt{id\_motivo} & UUID & FK $\rightarrow$ MOTIVO\_ CANCELACION\_ ENTREGA, NOT NULL & Motivo seleccionado del catálogo normalizado de causas de cancelación. \\
\addlinespace
\texttt{id\_estacion} & UUID & FK $\rightarrow$ ESTACION, NOT NULL & Estación de entrega desde donde se procesó la cancelación. \\
\addlinespace
\texttt{observaciones} & TEXT & NULL & Notas adicionales opcionalmente capturadas por el operador de entrega. \\
\addlinespace
\texttt{fecha} & TIMESTAMP WITH TIME ZONE & NOT NULL, DEFAULT NOW() & Marca temporal de la cancelación. \\
\end{tabularx}
\endgroup

\subsubsection{CORTE\_DE\_CAJA}
\begingroup
\setlength{\LTleft}{\fill}
\setlength{\LTright}{\fill}

\begin{tabularx}{\textwidth}{l p{2.5cm} p{3.5cm} X}
\caption{\textsc{Tabla CORTE\_DE\_CAJA de la base de datos.}} \label{tab:corte-de-caja}\\
\toprule
\textbf{Nombre del campo} & \textbf{Tipo de dato} & \textbf{Restricciones} & \textbf{Descripción funcional} \\
\midrule
\endfirsthead

\toprule
\textbf{Nombre del campo} & \textbf{Tipo de dato} & \textbf{Restricciones} & \textbf{Descripción funcional} \\
\midrule
\endhead

\midrule
\multicolumn{4}{r}{\small Continúa en la página siguiente...} \\
\endfoot
\bottomrule
\endlastfoot

\texttt{id\_corte} & UUID & PK, NOT NULL & Identificador único del corte de caja (UUID v4). \\
\addlinespace
\texttt{id\_estacion} & UUID & FK $\rightarrow$ ESTACION, NOT NULL & Estación de caja sobre la que se realiza el corte. \\
\addlinespace
\texttt{id\_usuario} & UUID & FK $\rightarrow$ USUARIO, NOT NULL & Usuario Manager que autorizó o ejecutó el corte. \\
\addlinespace
\texttt{monto\_inicial} & DECIMAL(10,2) & NOT NULL & Efectivo declarado en apertura de turno. \\
\addlinespace
\texttt{monto\_final\_declarado} & DECIMAL(10,2) & NULL & Efectivo contado físicamente al cierre. NULL mientras el corte no se finaliza. \\
\addlinespace
\texttt{monto\_calculado} & DECIMAL(10,2) & NULL & Total calculado por el sistema a partir de ventas en efectivo del turno. \\
\addlinespace
\texttt{diferencia} & DECIMAL(10,2) & NULL & Resultado de $\text{monto\_final\_declarado} - \text{monto\_calculado}$. Positivo = sobrante; negativo = faltante. \\
\addlinespace
\texttt{fecha\_apertura} & TIMESTAMP WITH TIME ZONE & NOT NULL & Marca temporal de inicio del turno. \\
\addlinespace
\texttt{fecha\_cierre} & TIMESTAMP WITH TIME ZONE & NULL & Marca temporal de cierre. NULL mientras el corte está abierto. \\
\end{tabularx}
\endgroup

\subsubsection{CONFIGURACION\_SISTEMA}
\begingroup
\setlength{\LTleft}{\fill}
\setlength{\LTright}{\fill}

\begin{tabularx}{\textwidth}{l p{2.5cm} p{3.5cm} X}
\caption{\textsc{Tabla CONFIGURACION\_SISTEMA de la base de datos.}} \label{tab:configuracion-sistema}\\
\toprule
\textbf{Nombre del campo} & \textbf{Tipo de dato} & \textbf{Restricciones} & \textbf{Descripción funcional} \\
\midrule
\endfirsthead

\toprule
\textbf{Nombre del campo} & \textbf{Tipo de dato} & \textbf{Restricciones} & \textbf{Descripción funcional} \\
\midrule
\endhead

\midrule
\multicolumn{4}{r}{\small Continúa en la página siguiente...} \\
\endfoot
\bottomrule
\endlastfoot

\texttt{id\_configuracion} & UUID & PK, NOT NULL & Identificador único del parámetro de configuración (UUID v4). \\
\addlinespace
\texttt{clave} & VARCHAR(100) & UNIQUE, NOT NULL & Identificador de la configuración (ej. \texttt{umbral\_urgencia\_rojo}, \texttt{nombre\_cafeteria}, \texttt{imagen\_cafeteria}). \\
\addlinespace
\texttt{valor} & TEXT & NOT NULL & Valor del parámetro serializado como texto. El campo \texttt{tipo\_dato} indica cómo deserializarlo. \\
\addlinespace
\texttt{tipo\_dato} & VARCHAR(20) & NOT NULL, CHECK (STRING $|$ INTEGER $|$ DECIMAL $|$ BOOLEAN $|$ TIME) & Tipo semántico del valor almacenado. Permite al cliente deserializar correctamente. \\
\addlinespace
\texttt{descripcion} & VARCHAR(255) & NULL & Explicación del propósito del parámetro para el administrador. \\
\end{tabularx}
\endgroup

\subsubsection{HISTORIAL\_DE\_MOVIMIENTOS}
\begingroup
\setlength{\LTleft}{\fill}
\setlength{\LTright}{\fill}

\begin{tabularx}{\textwidth}{l p{2.5cm} p{3.5cm} X}
\caption{\textsc{Tabla HISTORIAL\_DE\_MOVIMIENTOS de la base de datos.}} \label{tab:historial-de-movimientos}\\
\toprule
\textbf{Nombre del campo} & \textbf{Tipo de dato} & \textbf{Restricciones} & \textbf{Descripción funcional} \\
\midrule
\endfirsthead

\toprule
\textbf{Nombre del campo} & \textbf{Tipo de dato} & \textbf{Restricciones} & \textbf{Descripción funcional} \\
\midrule
\endhead

\midrule
\multicolumn{4}{r}{\small Continúa en la página siguiente...} \\
\endfoot
\bottomrule
\endlastfoot

\texttt{id\_historial} & UUID & PK, NOT NULL & Identificador único del evento auditado (UUID v4). \\
\addlinespace
\texttt{modulo} & VARCHAR(30) & NOT NULL & Módulo del sistema que generó el evento (ej. \texttt{MANAGER}, \texttt{CAJA}, \texttt{COCINA}, \texttt{ENTREGA}, \texttt{CLIENTE}). \\
\addlinespace
\texttt{accion} & VARCHAR(100) & NOT NULL & Descripción de la acción realizada (ej. \texttt{CREAR\_PRODUCTO}, \texttt{CAMBIAR\_ESTADO\_PEDIDO}). \\
\addlinespace
\texttt{id\_usuario} & UUID & FK $\rightarrow$ USUARIO, NULL & Usuario Manager o Cajero que ejecutó la acción. Nullable si el actor es un empleado o cliente. \\
\addlinespace
\texttt{id\_empleado} & UUID & FK $\rightarrow$ EMPLEADO, NULL & Empleado de cocina que ejecutó la acción. Nullable si el actor es un usuario o cliente. \\
\addlinespace
\texttt{id\_estacion} & UUID & FK $\rightarrow$ ESTACION, NULL & Estación desde la que se ejecutó la acción. Nullable si fue desde el módulo cliente. \\
\addlinespace
\texttt{id\_cliente} & UUID & FK $\rightarrow$ CLIENTE, NULL & Cliente que ejecutó la acción. Nullable si fue desde un puesto interno. \\
\addlinespace
\texttt{datos\_previos} & JSONB & NULL & Representación JSON del estado del registro antes del cambio. NULL para acciones de creación. \\
\addlinespace
\texttt{datos\_nuevos} & JSONB & NULL & Representación JSON del estado del registro después del cambio. NULL para acciones de eliminación. \\
\addlinespace
\texttt{fecha} & TIMESTAMP WITH TIME ZONE & NOT NULL, DEFAULT NOW() & Marca temporal exacta del evento auditado. \\
\end{tabularx}
\endgroup

\subsubsection{TABLAS DE UNIÓN Y DETALLES COMPLEMENTARIAS}
Las siguientes tablas completan el modelo; su estructura es sencilla y se describe brevemente:
\begingroup
\setlength{\LTleft}{\fill}
\setlength{\LTright}{\fill}

\begin{tabularx}{\textwidth}{l p{6cm} X}
\caption{\textsc{Tablas de unión y detalles complementarias de la base de datos.}} \label{tab:tablas-union-detalles}\\
\toprule
\textbf{Tabla} & \textbf{Campos} & \textbf{Propósito} \\
\midrule
\endfirsthead

\toprule
\textbf{Tabla} & \textbf{Campos} & \textbf{Propósito} \\
\midrule
\endhead

\midrule
\multicolumn{3}{r}{\small Continúa en la página siguiente...} \\
\endfoot
\bottomrule
\endlastfoot

\texttt{ESTACION\_ETIQUETA} & \texttt{id\_estacion} FK, \texttt{id\_etiqueta} FK (PK compuesta) & Relación M:N entre estaciones y etiquetas de clasificación. \\
\addlinespace
\texttt{INGREDIENTE\_PRODUCTO} & \texttt{id\_ingrediente\_producto} PK, \texttt{id\_producto} FK, \texttt{id\_insumo} FK, \texttt{cantidad\_default}, \texttt{es\_personalizable}, \texttt{es\_opcional} & Define la composición base de cada producto para personalización y descuento de inventario. \\
\addlinespace
\texttt{INSUMO\_DESECHABLE\_PRODUCTO} & \texttt{id\_insumo\_desechable} PK, \texttt{id\_producto} FK, \texttt{id\_insumo} FK, \texttt{cantidad} & Registra materiales no alimentarios asociados a un producto (vasos, bolsas, servilletas). \\
\addlinespace
\texttt{PRODUCTO\_PROMOCION} & \texttt{id\_producto} FK, \texttt{id\_promocion} FK (PK compuesta) & Relación M:N que vincula productos con las promociones aplicables. \\
\addlinespace
\texttt{DEPENDENCIA\_PASO} & \texttt{id\_paso\_receta} FK, \texttt{id\_paso\_receta\_previo} FK (PK compuesta) & Aristas del DAG de dependencias entre pasos de receta del catálogo. \\
\addlinespace
\texttt{INGREDIENTE\_PASO} & \texttt{id\_ingrediente\_paso} PK, \texttt{id\_paso\_receta} FK, \texttt{id\_insumo} FK, \texttt{cantidad} & Insumos consumidos en cada paso de receta del catálogo. \\
\addlinespace
\texttt{DEPENDENCIA\_PASO\_PEDIDO} & \texttt{id\_paso\_receta\_pedido} FK, \texttt{id\_paso\_receta\_pedido\_previo} FK (PK compuesta) & Instantánea inmutable del DAG de dependencias para el pedido en curso. \\
\addlinespace
\texttt{INGREDIENTE\_PASO\_PEDIDO} & \texttt{id\_ingrediente\_paso\_pedido} PK, \texttt{id\_paso\_receta\_pedido} FK, \texttt{id\_insumo} FK, \texttt{cantidad\_snapshot} & Insumos por consumir en cada paso de producción, congelados al confirmar el pedido. \\
\addlinespace
\texttt{AJUSTE\_CORTE\_CAJA} & \texttt{id\_ajuste} PK, \texttt{id\_corte} FK, \texttt{metodo\_pago}, \texttt{monto}, \texttt{descripcion} & Desglose de ventas por método de pago dentro de un corte de caja. \\
\addlinespace
\texttt{AJUSTE\_CORTE\_INVENTARIO} & \texttt{id\_ajuste\_inventario} PK, \texttt{id\_corte} FK, \texttt{id\_insumo} FK, \texttt{cantidad\_contada}, \texttt{cantidad\_sistema}, \texttt{diferencia} & Conciliación de inventario físico contra sistema para cada insumo al cierre de turno. \\
\addlinespace
\texttt{MOTIVO\_CANCELACION\_ENTREGA} & \texttt{id\_motivo} PK, \texttt{descripcion}, \texttt{activo} & Catálogo normalizado de causas de cancelación seleccionables en el módulo de Entrega. \\
\addlinespace
\texttt{CONFIGURACION\_METODO\_PAGO} & \texttt{id\_metodo\_pago} PK, \texttt{nombre} UNIQUE, \texttt{habilitado}, \texttt{descripcion} & Métodos de pago disponibles en el sistema, habilitables individualmente por el Manager. \\
\end{tabularx}
\endgroup

\newpage
\section{DIAGRAMAS DE SECUENCIA}
Los diagramas de secuencia integrados en la presente documentación modelan la dimensión cronológica de los intercambios de mensajes entre los participantes de los procesos críticos del sistema. Bajo la notación UML 2.5, las líneas sólidas representan los mensajes síncronos, mientras que las líneas discontinuas simbolizan las respuestas correspondientes. La estructura de los bloques \textit{alt/else} representa los flujos alternativos y mutuamente excluyentes, derivados de las condiciones de negocio definidas, y los bloques \textit{loop} denotan la iteración sobre colecciones de entidades. La arquitectura de los participantes se segmenta en las capas definidas: capa de presentación (Frontend Next.js), capa de aplicación (API Gateway y servicios Express.js), capa de persistencia (PostgreSQL) y el canal de comunicación en tiempo real (WebSocket/SSE).

\subsection{DS-01: CONFIRMACIÓN DE PEDIDO DESDE EL MÓDULO CLIENTE}
El proceso de confirmación de pedido constituyó el flujo transaccional de mayor complejidad dentro del sistema, dado que involucró cuatro capas de validación secuencial (RN-CLI-018), la aplicación del patrón de inmutabilidad mediante instantáneas sobre precios e ingredientes (RN-MGR-021 y RN-MGR-023), el control de condiciones de carrera sobre la capacidad del slot de recolección (RN-CLI-017, RN-CAJ-021) y la inicialización del grafo de estados del KDS al término de la transacción.

El diagrama expuso tres características de diseño determinantes. En primer lugar, la validación del token de sesión operó como guarda transversal en el API Gateway antes de que se invocara cualquier lógica de negocio, lo cual materializó el requisito RNFU-01 relativo al control de acceso por rol. En segundo lugar, las validaciones de existencias y capacidad de slot se ejecutaron de forma secuencial antes del inicio de la transacción de escritura. Este orden no resultó arbitrario, pues la consulta de existencias no requirió bloqueo exclusivo, mientras que la consulta de capacidad de slot emitió una sentencia SQL \texttt{SELECT FOR UPDATE}, la cual previno la condición de carrera identificada en los requisitos RN-CLI-017 y RN-CAJ-021. Finalmente, todas las operaciones de inserción que constituyeron el pedido, incluyendo las instantáneas de precios, la receta y la generación del token QR, se envolvieron en una única transacción atómica. En consecuencia, cualquier fallo en el proceso provocó un retroceso (\textit{rollback}) completo sin estados parciales persistidos, con lo cual se satisfizo el requisito RNFU-04.

\begin{landscape}
\begin{figure}[H]
    \centering
    \includesvg[width=1.35\textwidth, inkscapelatex=false]{diagrams/DA-01/DA-01 Ciclo de Vida del Pedido.svg}
    \caption{Diagrama de secuencia confirmación de pedido.}
    \label{fig:diagrama-secuencia-confirmacion-pedido}
\end{figure}
\end{landscape}

\subsection{DS-02: COMPLETAR PASO DE PREPARACIÓN Y EVALUACIÓN DEL GRAFO DE DEPENDENCIAS (DAG)}
El presente proceso modela la transición de un paso al estado \texttt{COMPLETADO} dentro del módulo Cocina, operación que conlleva la evaluación recursiva del grafo de dependencias (DAG) para determinar la liberación de los pasos subsiguientes. Este flujo constituye la acción con mayor impacto transversal del sistema, debido a que su ejecución puede desencadenar la transición del pedido al estado \texttt{PREPARADO} y propagar notificaciones simultáneas hacia el módulo Entrega y el módulo Cliente.

El diseño del proceso se sustentó en dos invariantes fundamentales. En primer lugar, se estableció una separación estricta de responsabilidades entre la transición de estado del paso, gestionada por \texttt{StepService}, y la evaluación del grafo de dependencias, a cargo de \texttt{DAGService}. Esta arquitectura permitió que el componente de evaluación operara como una entidad autónoma y reutilizable, lo cual facilitó su invocación durante la inicialización del pedido en el módulo Caja. En segundo lugar, se implementó la propagación selectiva de eventos mediante un canal de comunicación en tiempo real. Esta solución posibilitó la liberación de tareas en múltiples estaciones del KDS, el cambio de estado del pedido a \texttt{PREPARADO} y la actualización inmediata de la interfaz del usuario sin necesidad de sondeo (\textit{polling}) o recarga manual, con lo cual se satisfizo el requisito REQ-GLB-001 (RN-SYS-003).
\begin{figure}[H]
    \centering
    \includesvg[width=0.55\textwidth, inkscapelatex=false]{diagrams/DA-02/DA-02 Flujo de Confirmación de Pedido con Validaciones Transaccionales.svg}
    \caption{Diagrama de secuencia de completar paso de preparación.}
    \label{fig:diagrama-secuencia-completar-paso}
\end{figure}
\subsection{DS-03: PROTOCOLO DE VERIFICACIÓN QR Y CONFIRMACIÓN DE ENTREGA}
El protocolo de verificación mediante código QR constituyó la fase terminal del ciclo de vida del pedido. Este proceso requirió la coordinación de tres participantes distintos —el cliente, el operador de entrega y el hardware periférico— y la ejecución de cuatro condiciones de validación secuencial sobre el token (RN-ENT-014). El flujo transaccional culminó con la escritura atómica del estado \texttt{ENTREGADO} en la base de datos, operación que precedió obligatoriamente al renderizado de la confirmación visual (RN-ENT-017).

El diseño del protocolo operó de manera simétrica para clientes autenticados (CLI-14) y anónimos (CLI-16). El servidor validó exclusivamente los atributos del \texttt{CodigoVerificacionQR.token} recibido, omitiendo cualquier distinción sobre el origen de la solicitud de generación o la necesidad de un token de sesión para procesar la validación dentro del módulo Entrega (RN-CLI-029). Esta estrategia de diseño garantizó la consistencia operativa independientemente del estado de autenticación del usuario.

Respecto al ciclo de vida del token, se implementó una política de invalidación permanente ante cualquier condición de fallo, incluida la expiración del tiempo de vida (TTL). El token cambió invariablemente a un estado de \texttt{CONSUMIDO} o \texttt{EXPIRADO}, lo cual impidió su reutilización bajo cualquier circunstancia (RN-ENT-016). En consecuencia, esta restricción forzó al usuario a generar un nuevo código QR desde el módulo cliente (CLI-14) cuando el proceso de recolección requirió un reintento.

\begin{landscape}
\begin{figure}[H]
    \centering
    \includesvg[width=1.25\textwidth, inkscapelatex=false]{diagrams/DA-03/DA-03 Evaluación del DAG y Propagación de Estado en el KDS.svg}
    \caption{Diagrama de secuencia de protocolo de verificación de QR.}
    \label{fig:diagrama-secuencia-verificacion-qr}
\end{figure}
\end{landscape}

\subsection{DS-04: CREACIÓN DE PEDIDO PRESENCIAL EN MÓDULO CAJA CON TICKET DIGITAL}
El flujo de creación de pedido en el Módulo Caja constituyó el segundo proceso transaccional de mayor complejidad dentro del sistema, cuya criticidad radicó en la gestión de dos tipos de pedido con bifurcaciones operativas distintas: \texttt{INMEDIATO} y \texttt{ANTICIPADO}. El diseño de este proceso se caracterizó por cuatro pilares estructurales: la ejecución del decremento atómico del inventario en el momento de la confirmación del pago (RN-CAJ-009), la generación condicional de un token de acceso anónimo para pedidos sin cuenta asociada (RN-CLI-025) ---el cual sustituyó integralmente al comprobante físico--- y la validación de capacidad de ventana mediante concurrencia con \texttt{SELECT FOR UPDATE} en pedidos anticipados (RN-CAJ-021). 

El diagrama expuso una asimetría fundamental en el modelo de autenticación respecto al módulo Cliente. Mientras que este último operó bajo credenciales individuales por cuenta de usuario, el Módulo Caja autenticó mediante credenciales compartidas de estación (\texttt{CredencialEstacion}), con lo cual se vinculó la trazabilidad operativa al \texttt{id\_estacion} en lugar de a un individuo específico (RN-CAJ-002). 

En cuanto a la integridad de los datos, el decremento de existencias se ejecutó dentro de la misma transacción atómica que creó el pedido. Esta arquitectura eliminó la existencia de estados intermedios inconsistentes en los cuales el sistema registraba el pedido sin haber procesado la reducción correspondiente del inventario (RN-CAJ-009). Asimismo, la generación del token de acceso anónimo operó bajo una lógica condicional; el servidor únicamente creó el registro en \texttt{PedidoAccesoAnonimo} cuando el atributo \texttt{Pedido.cliente\_id} resultó nulo. Este diseño transformó el ticket impreso en una URL de seguimiento digital, cumpliendo así con el requisito RN-CLI-025 sin imponer una carga administrativa de registro de cuentas. 

Finalmente, la asignación de ventana en los pedidos anticipados funcionó como una transacción aislada que emitió su propio \texttt{SELECT FOR UPDATE}. Esta segmentación desacopló la resolución de la condición de carrera de la creación general del pedido. Como resultado, el sistema permitió que el modal de selección permaneciera activo para reintentos en caso de que la capacidad se agotara durante el intervalo entre la carga de la interfaz y la confirmación del usuario (RN-CAJ-021). 
\begin{figure}[H]
    \centering
    \includesvg[width=0.75\textwidth, inkscapelatex=false]{diagrams/DA-04/DA-04 Ciclo Completo del Corte de Caja con Doble Reconciliación.svg}
    \caption{Diagrama de secuencia de creación de pedido presencial.}
    \label{fig:diagrama-secuencia-creacion-pedido-presencial}
\end{figure}
\subsection{DS-05: ALTA DE EMPLEADO Y FLUJO DE ONBOARDING}
El proceso de alta de un nuevo empleado constituyó un flujo de trabajo multi-actor que involucró la intervención del Manager, el servidor SMTP como sistema externo de mensajería y la sesión del empleado durante la ejecución de su activación mediante un enlace de uso único. El diseño del sistema distinguió el modelo de autenticación del KDS, el cual requirió la identificación individual del personal, frente al modelo de credenciales compartidas empleado en los módulos de Caja y Entrega. Debido a esta arquitectura, el proceso de onboarding determinó la capacidad operativa del entorno de cocina. En cumplimiento con los requisitos RN-MGR-007 y RN-MGR-010, el sistema ejecutó validaciones estrictas sobre la compatibilidad del tipo de estación y la unicidad del correo electrónico antes de la persistencia de cualquier entidad. Asimismo, la asignación del valor \texttt{PENDIENTE\_CONTRASENA} al campo \texttt{Empleado.estado} garantizó, bajo el cumplimiento de RN-MGR-009, que ningún empleado autenticara su acceso sin haber completado previamente la secuencia de activación. 

El diseño del flujo integró tres propiedades fundamentales derivadas de los requerimientos de seguridad y operación. En primer lugar, la arquitectura segregó la creación del empleado de la asignación de estación en dos transacciones independientes. Esta separación se implementó debido a que la validación de la incompatibilidad del tipo de estación (RN-MGR-007) requirió, como prerrequisito funcional, la disponibilidad del identificador del empleado (\texttt{id\_empleado}) recién generado. En segundo lugar, la comunicación con el servidor SMTP se orquestó fuera del alcance de la transacción de la base de datos. Esta estrategia aseguró que la operación de \texttt{COMMIT} precediera al envío del correo electrónico; en caso de producirse un fallo en el servicio de mensajería, el registro del empleado permaneció persistido, lo cual permitió al Manager solicitar el reenvío de la notificación sin incurrir en la duplicidad del alta. 

En tercer lugar, el mecanismo de activación empleó un token opaco en el enlace de onboarding en lugar de exponer el identificador del empleado en la URL. Esta implementación siguió los principios de opacidad definidos en RN-CLI-027, con lo cual se mitigó eficazmente el riesgo de enumeración de identificadores de empleados mediante la ejecución de peticiones consecutivas hacia la ruta de activación. 

\begin{figure}[H]
    \centering
    \includesvg[width=0.7\textwidth, inkscapelatex=false]{diagrams/DA-05/DA-05 Registro de Producto con Validación del Invariante de Distribución de Receta.svg}
    \caption{Diagrama de secuencia de alta de empleado.}
    \label{fig:diagrama-secuencia-alta-empleado}
\end{figure}

\subsection{DS-06: CORTE DE CAJA Y CIERRE DE TURNO}
El proceso de corte de caja representó el punto de mayor complejidad en cuanto a la validación de la integridad financiera dentro del sistema. Su criticidad operativa residió en la implementación de un patrón de resolución iterativa, mediante el cual cada inconsistencia ---ya fuera de índole financiera o de inventario--- impidió el avance del proceso hasta su justificación explícita por parte del cajero, sin posibilidad de omisión. Se aplicó la fórmula de (RN-CAJ-016), la cual exigió al usuario justificar cada discrepancia mediante el registro de una entidad \texttt{AjusteCorteCaja}. De forma análoga, el control de inventario bloqueó el diálogo de confirmación de cierre ante la existencia de discrepancias entre el stock del sistema y el conteo físico, conforme a lo establecido en la restricción RN-CAJ-017. La confirmación final del cierre constituyó una operación irreversible que invalidó la sesión del usuario de manera atómica con el registro del evento de auditoría correspondiente (RN-CAJ-018). 

El diseño del proceso integró tres propiedades técnicas distintivas que lo diferenciaron del resto de los flujos del sistema. En primer lugar, las fases de reconciliación financiera e inventario se ejecutaron de forma secuencial y no paralela; el sistema denegó el inicio del conteo de inventario hasta que las discrepancias de caja fueron subsanadas, lo cual garantizó la ausencia de inconsistencias silenciosas en ambas dimensiones. En segundo lugar, la invalidación de la sesión del cajero (RN-CAJ-018) se orquestó en el servidor tras la confirmación exitosa de la transacción de cierre. Si bien un fallo en la operación de \texttt{COMMIT} permitió la persistencia de la sesión para reintentos, el éxito de la transacción tornó la sesión irrecuperable, lo que aseguró la trazabilidad definitiva del evento. En tercer lugar, el resumen del turno incluyó de forma explícita los pedidos con estado \texttt{NO\_PAGADO}, lo cual permitió al cajero regularizar deudas pendientes mediante la interfaz de gestión de pagos (RN-CAJ-014) antes de formalizar el cierre, previniendo así la existencia de ingresos no registrados en el balance final del \texttt{CorteDeCaja}. 

\begin{landscape}
\begin{figure}[H]
    \centering
    \includesvg[width=1.35\textwidth, inkscapelatex=false]{diagrams/DA-06/DA-06 Algoritmo de Cálculo de Urgencia Dual en KDS y Módulo Entrega.svg}
    \caption{Diagrama de secuencia de corte de caja y cierre de turno.}
    \label{fig:diagrama-secuencia-corte-caja}
\end{figure}
\end{landscape}

\subsection{DS-07: CANCELACIÓN DE PEDIDO DESDE ENTREGA CON PROPAGACIÓN ASÍNCRONA}
El proceso de cancelación de un pedido desde el módulo Entrega constituyó el único flujo del sistema capaz de generar efectos de escritura en cascada sobre múltiples módulos de forma sincrónica dentro de una misma transacción, además de propagar notificaciones asíncronas diferenciadas según el estado del pedido al momento de la ejecución. La criticidad de este proceso residió en tres invariantes estructurales. En primer lugar, la justificación de la cancelación mediante un motivo obligatorio, normalizado contra el catálogo \texttt{MotivoCancelacionEntrega} previo a la persistencia, garantizó la trazabilidad operativa (RN-ENT-006). En segundo lugar, la captura del estado del pedido (\texttt{estado\_al\_cancelar}) se ejecutó dentro de la misma transacción mediante un bloqueo \texttt{FOR UPDATE}, operación indispensable para determinar la invalidación de los pasos en el KDS (RN-ENT-008, RN-ENT-009). Finalmente, la transición al estado \texttt{CANCELADO} operó como una condición terminal e irreversible, restringiendo cualquier modificación posterior desde cualquier interfaz del sistema (RN-ENT-007). 

El diseño del diagrama expuso tres propiedades de alta consecuencia técnica. La primera fue la bifurcación de efectos secundarios determinada por el estado del pedido: una cancelación efectuada desde \texttt{EN\_PREPARACION} invalidó todos los registros \texttt{PasoPedidoEstado} con estatus \texttt{PENDIENTE} o \texttt{BLOQUEADO} ---ocasionando su remoción inmediata del display del KDS---, mientras que una cancelación desde \texttt{PREPARADO} no produjo efecto sobre el KDS, dado que los pasos completados se trataron como registros históricos inmutables (RN-ENT-008). La segunda propiedad concernió a la normalización del catálogo de motivos, la cual se ejecutó fuera del bloque \texttt{BEGIN TRANSACTION}; esta estrategia evitó que un bloqueo de fila en la tabla de motivos retrasara innecesariamente la cancelación del pedido (RN-ENT-006). La tercera fue la propagación condicional de eventos hacia el módulo Cliente: el servidor emitió el evento al Frontend Cliente exclusivamente cuando el atributo \texttt{Pedido.cliente\_id} poseyó un valor no nulo, lo cual evitó el procesamiento innecesario de notificaciones para pedidos anónimos originados desde el módulo Caja que carecían de una sesión de usuario activa. 

\begin{figure}[H]
    \centering
    \includesvg[width=1\textwidth, inkscapelatex=false]{diagrams/DA-07/DA-07 Flujo de Cancelación con Propagación Diferenciada por Estado del Pedido.svg}
    \caption{Diagrama de secuencia de cancelación de pedido en la entrega.}
    \label{fig:diagrama-secuencia-cancelacion-pedido}
\end{figure}    

\section{DISEÑO DE INTERFACES}
El prototipo de interfaces fue desarrollado en Figma y contempla los siguientes flujos principales. 

\subsection{FLUJO DEL CLIENTE}
El cliente accede mediante autenticación con su correo institucional. La pantalla principal muestra el menú disponible organizado por categorías, con indicadores visuales de disponibilidad. El usuario selecciona sus productos, los agrega al carrito, elige el horario de entrega disponible dentro del turno activo, y confirma el pedido. Una pantalla de seguimiento permite consultar el estado en tiempo real.

\begin{figure}[H]
    \centering
    \includesvg[width=0.4\textwidth]{images/cliente/CLI-01 — Cliente-Login.svg}
    \caption{Interfaz de inicio de sesión para el cliente.}
    \label{fig:ui-cli-01}
\end{figure}

Esta interfaz corresponde a la pantalla de inicio de sesión. Presenta una estructura minimalista que incluye un formulario de autenticación con dos campos de captura principales: uno para ingresar el correo electrónico del usuario y otro para la contraseña asociada. En la parte inferior del bloque central se dispone de un botón de acción principal etiquetado como ``Iniciar sesión'' para procesar las credenciales, complementado con dos accesos directos en la base de la pantalla: uno para la recuperación de cuentas mediante la opción ``¿Olvidaste tu contraseña?'' y otro para redirigir a nuevos usuarios a través del enlace ``Registrarse''.

\begin{figure}[H]
    \centering
    \includesvg[width=0.4\textwidth]{images/cliente/CLI-02 — Cliente-Home.svg}
    \caption{Interfaz principal para el cliente.}
    \label{fig:ui-cli-02}
\end{figure}

Esta interfaz corresponde a la pantalla principal o tablero de inicio. El diseño se estructura mediante un menú de navegación superior que permite al usuario explorar el catálogo de productos a través de diferentes categorías comerciales, tales como ``Paquetes'', ``Bebidas'', ``Comida'', ``Snacks'' y ``Postres''. En la sección central se despliegan de forma tabular las opciones destacadas o ``Paquetes del día'', donde cada artículo se presenta en tarjetas visuales independientes que detallan el nombre del combo, los componentes específicos de alimentos y bebidas que lo integran, su costo unitario y un botón interactivo para agregar el elemento al carrito de compra de manera directa.

\begin{figure}[H]
    \centering
    \includesvg[width=0.4\textwidth]{images/cliente/CLI-03 — Cliente-BusquedaResultados.svg}
    \caption{Interfaz de búsqueda para el cliente.}
    \label{fig:ui-cli-03}
\end{figure}

Esta interfaz ilustra la funcionalidad del motor de búsqueda integrado en el módulo del cliente. Al interactuar con el control de exploración, la pantalla despliega una barra de búsqueda en la sección superior para la introducción de texto libre, complementada por un historial de ``Búsquedas recientes'' en la parte inferior para agilizar la navegación. El cuerpo principal de la interfaz genera un listado dinámico con los resultados que coinciden con el criterio ingresado, organizando los productos disponibles en tarjetas individuales que muestran la fotografía del artículo, su nombre comercial, el precio unitario y un botón de acción rápida para agregar el producto directamente al carrito de compra.

\begin{figure}[H]
    \centering
    \includesvg[width=0.4\textwidth]{images/cliente/CLI-04 — Cliente-Promociones.svg}
    \caption{Interfaz de catálogo de promociones para el cliente.}
    \label{fig:ui-cli-04}
\end{figure}

Esta interfaz corresponde a la vista del catálogo filtrada específicamente para la sección de ``Promociones'' dentro del módulo del cliente. La pantalla mantiene la consistencia de navegación superior y el banner informativo de disponibilidad horaria, organizando los productos en una cuadrícula o \textit{grid} de tarjetas individuales dispuestas en tres columnas. Cada tarjeta detalla la imagen del platillo o combo, su nombre comercial, el precio correspondiente y un indicador de disponibilidad en la base; este último cambia de estado y color (como el ``Combo Asados'' marcado en rojo como ``No disponible'') para informar en tiempo real al usuario si el artículo puede ser añadido al carrito de compras en ese momento.

\begin{figure}[H]
    \centering
    \includesvg[width=0.4\textwidth]{images/cliente/CLI-05 — Cliente-Productos.svg}
    \caption{Interfaz de detalles de la promoción seleccionada.}
    \label{fig:ui-cli-05}
\end{figure}

Esta interfaz corresponde a la vista de detalle de un producto seleccionado, dentro del módulo del cliente. En la sección superior se incluye la opción de navegación ``Regresar'', acompañada por el nombre comercial de la promoción, su costo total y una breve descripción general de los artículos incluidos. El cuerpo principal de la pantalla desglosa el apartado ``Contenido'', organizando en tarjetas independientes cada uno de los elementos que integran el paquete, detallando sus ingredientes o método de elaboración junto con una imagen de referencia.

En la parte inferior de la interfaz se integra una barra de herramientas operativa para la gestión del pedido. Esta sección incluye un selector numérico flanqueado por controles interactivos de incremento (``+'') y decremento (``-'') para parametrizar la cantidad de unidades deseadas. Asimismo, se disponen dos botones de acción principal en paralelo: ``Agregar a la orden'', para consolidar los artículos en el carrito de compras, y ``Quitar de la orden'', resaltado en color rojo para revertir la selección.

\begin{figure}[H]
    \centering
    \includesvg[width=0.4\textwidth]{images/cliente/CLI-06 — Cliente-DetalleCombo.svg}
    \caption{Interfaz de catálogo de productos para el cliente.}
    \label{fig:ui-cli-06}
\end{figure}

Esta interfaz corresponde a la vista del catálogo filtrada para la sección general de ``Productos'' dentro del módulo del cliente. El diseño conserva la estructura de navegación superior, la barra de búsqueda y el banner informativo que indica la disponibilidad de espacios por bloques horarios.

Los artículos se organizan de manera simétrica en un formato de rejilla de tres columnas, donde cada tarjeta individual muestra la fotografía del platillo, su nombre comercial y el costo unitario. En la parte inferior de cada bloque se integra un validador de existencias en tiempo real; los productos listos para su adquisición se marcan con una verificación verde y la etiqueta ``Disponible'', mientras que aquellos agotados o fuera de inventario temporal se presentan atenuados en escala de grises con un indicador visual restrictivo en color rojo que cita ``No disponible'', inhabilitando su selección inmediata.

\begin{figure}[H]
    \centering
    \includesvg[width=0.4\textwidth]{images/cliente/CLI-07 — Cliente-DetalleComboIngredientes.svg}
    \caption{Interfaz de personalización de producto para el cliente.}
    \label{fig:ui-cli-07}
\end{figure}

Esta interfaz ilustra la vista detallada de los componentes e ingredientes de un producto y su parametrización de personalización. En la sección superior se observa el desglose del apartado ``Contenido'', el cual especifica de forma tabular los insumos base obligatorios de la receta estándar, detallando el nombre del ingrediente, la cantidad exacta y su unidad de medida correspondiente.

Inmediatamente abajo se despliega la sección ``Extras'', un componente interactivo que permite al consumidor personalizar su pedido agregando aditamentos opcionales mediante casillas de verificación. Al seleccionar una variante (como ``Cebolla Caramelizada'' o ``Salsa BBQ''), el sistema habilita un selector numérico para ajustar de manera exacta la porción en gramos o mililitros, calculando e indicando simultáneamente el costo adicional que se sumará al importe de la orden antes de su confirmación final en la base de la pantalla.

\begin{figure}[H]
    \centering
    \includesvg[width=0.4\textwidth]{images/cliente/CLI-08 — Cliente-Carrito.svg}
    \caption{Interfaz de carrito de compras para el cliente.}
    \label{fig:ui-cli-08}
\end{figure}

Esta interfaz corresponde a la vista del ``Carrito de compras'' dentro del módulo del cliente. En la parte superior se ubica un control de navegación con la etiqueta ``Regresar''. El área central contiene una lista de tarjetas para los productos que el usuario ha seleccionado. Cada una de estas tarjetas incluye la fotografía del artículo correspondiente y controles interactivos con un selector numérico para modificar las cantidades mediante los botones ``+'', ``-'' o el icono de eliminación.

En el extremo inferior se organiza la sección de liquidación, la cual presenta la etiqueta ``Total'' alineada al importe acumulado. Debajo de este balance se ubican dos botones de acción principal en paralelo: ``Vaciar carrito'', diseñado en color rojo, y ``Proceder al pago'', diseñado en color negro con un icono de tarjeta bancaria.

\begin{figure}[H]
    \centering
    \includesvg[width=0.4\textwidth]{images/cliente/CLI-09 — Cliente-AvisoContraEntrega.svg}
    \caption{Interfaz de validación de condiciones de entrega.}
    \label{fig:ui-cli-09}
\end{figure}

Esta interfaz corresponde a la pantalla de validación de condiciones de entrega dentro del flujo del carrito de compras del módulo del cliente. En la sección superior se sitúa el control de navegación "Regresar". El cuerpo central muestra un aviso con el encabezado "Recuerda que solo se aceptan pagos a contra entrega", seguido de un párrafo informativo sobre el compromiso con la comunidad, la puntualidad en la recolección y la prevención de penalizaciones. Justo debajo, se integran dos botones de opción interactivos para que el usuario elija entre las alternativas de "Aceptar" o "Rechazar" los términos expuestos. 

En el extremo inferior se organiza la sección de liquidación, la cual expone la etiqueta "Total" alineada a un importe destacado. En la base se disponen dos botones de acción en paralelo: "Cancelar", diseñado con un icono de cruz sobre un fondo de color rojo, y "Seleccionar recolección", estructurado con un icono de calendario sobre un fondo negro. 

\begin{figure}[H]
    \centering
    \includesvg[width=0.4\textwidth]{images/cliente/CLI-10 — Cliente-SeleccionHorario.svg}
    \caption{Interfaz de selección de horarios para el cliente.}
    \label{fig:ui-cli-10}
\end{figure}

En la sección superior se sitúa el control de navegación ``Regresar''. El cuerpo central presenta el encabezado ``Seleccione un horario para recolectar su pedido'', seguido de un listado vertical de tarjetas que representan los bloques de tiempo disponibles. Cada tarjeta muestra un rango horario específico y un indicador que detalla la cantidad de espacios disponibles en tiempo real, complementado con un icono de flecha para la selección del bloque.

En el extremo inferior se organiza la sección de cierre, la cual expone la etiqueta ``Total'' junto al balance acumulado. En la base se disponen dos botones de acción en paralelo: ``Cancelar'', diseñado con un icono de cruz sobre un fondo de color rojo, y ``Ver el resumen'', estructurado con un icono de documento sobre un fondo negro.

\begin{figure}[H]
    \centering
    \includesvg[width=0.4\textwidth]{images/cliente/CLI-11 — Cliente-ResumenPedido.svg}
    \caption{Interfaz de confirmación de resumen del pedido para el cliente.}
    \label{fig:ui-cli-11}
\end{figure}

Esta interfaz corresponde a la pantalla de ``Resumen del pedido'' dentro del flujo final del carrito de compras. En la parte superior se sitúa el control de navegación ``Regresar''. El cuerpo principal organiza la información fija del proceso comenzando con el apartado de ``Contenido'', el cual despliega una lista vertical con las tarjetas de los artículos adquiridos, especificando para cada uno su fotografía, nombre, descripción o aditamentos seleccionados y la cantidad de unidades solicitadas en la base de cada tarjeta. Inmediatamente abajo se encuentra la sección de ``Método de pago'', la cual indica de forma textual la forma de liquidación seleccionada por el consumidor detallada como ``Tarjeta de débito'', seguida por el apartado de ``Horario seleccionado'' que muestra el bloque de tiempo previamente reservado por el usuario para acudir por su orden.

En el extremo inferior se organiza la sección de cierre de la transacción, donde se presenta la etiqueta ``Total'' junto al balance acumulado. En la base de la interfaz se disponen dos botones de acción en paralelo, identificados como ``Cancelar'', diseñado con un icono de cruz sobre un fondo de color rojo, y ``Finalizar pedido'', estructurado con un icono de verificación sobre un fondo verde para consolidar el proceso en el sistema.

\begin{figure}[H]
    \centering
    \includesvg[width=0.4\textwidth]{images/cliente/CLI-12 — Cliente-ConfirmacionPedido.svg}
    \caption{Interfaz de confirmación y comprobante de pedido.}
    \label{fig:ui-cli-12}
\end{figure}

Esta interfaz corresponde a la pantalla de confirmación o comprobante tras procesar la orden. En la sección superior se sitúa un botón de navegación con la inscripción ``Regresar''. Inmediatamente abajo se despliega un mensaje de éxito que cita ``¡Pedido generado con éxito!'', acompañado por un identificador numérico de orden en tipografía de gran tamaño y el apartado ``Realizado el'' que detalla de forma textual la fecha y la hora del registro.

El cuerpo principal organiza la información fija de la transacción comenzando con el bloque de ``Contenido'', el cual muestra una lista vertical con las tarjetas de los artículos seleccionados, especificando para cada uno su fotografía, nombre comercial, descripción o aditamentos seleccionados y la cantidad de unidades en la base de cada tarjeta. Debajo se encuentra la sección de ``Método de pago'' que indica la alternativa de liquidación registrada como ``Tarjeta de débito'', seguida por el apartado de ``Horario seleccionado'' que muestra el bloque de tiempo reservado para acudir por el pedido. En el extremo inferior se presenta la etiqueta ``Total'' junto al balance acumulado, concluyendo la interfaz con un botón de acción principal de ancho completo que cita ``Continuar'', diseñado con texto blanco sobre un fondo negro.

\begin{figure}[H]
    \centering
    \includesvg[width=0.4\textwidth]{images/cliente/CLI-13 — Cliente-PedidosPorRecolectar.svg}
    \caption{Interfaz de recolección de pedido.}
    \label{fig:ui-cli-13}
\end{figure}

Esta interfaz corresponde a la pantalla de seguimiento denominada "Pedidos por recolectar". En el extremo superior izquierdo se dispone el control de navegación con la inscripción "Regresar". El área central organiza de manera vertical las tarjetas informativas de los pedidos activos en el sistema, identificados con los folios numéricos. Cada bloque detalla de forma textual la fecha y hora exacta del registro bajo la etiqueta "Realizado el", acompañada por una barra de progreso horizontal en color azul y una pestaña inferior negra interactiva con la palabra "Recolectar" y un icono de flecha. En la base de la pantalla se mantiene fija la barra de navegación global con las secciones de "Menú", "Recolección" y "Perfil", resaltando la pestaña central en color negro. 

\begin{figure}[H]
    \centering
    \includesvg[width=0.4\textwidth]{images/cliente/CLI-14 — Cliente-DetallePedidoRecoleccion.svg}
    \caption{Interfaz de recolección y resumen de pedido.}
    \label{fig:ui-cli-14}
\end{figure}

Esta interfaz corresponde a la pantalla de detalle y validación de una orden específica dentro del apartado de recolección. En el extremo superior izquierdo se dispone el control de navegación ``Regresar''. Inmediatamente abajo se expone el número de folio del pedido, seguido de un botón de acción principal que cita ``Escanear código'', el cual está estructurado con un icono de código QR sobre un fondo negro. Debajo se añade el indicador ``Realizado el'' que detalla de forma textual la fecha y hora del registro de la orden.

El resto del cuerpo de la pantalla organiza la información de la transacción comenzando con el bloque de ``Contenido'', el cual muestra una lista vertical con las tarjetas de los artículos solicitados especificando para cada uno su fotografía, nombre comercial, descripción y la cantidad de unidades. Posteriormente se encuentra la sección de ``Método de pago'' que indica la alternativa de liquidación registrada como ``Tarjeta de débito'', seguida por el apartado de ``Horario seleccionado''. En el extremo inferior de la interfaz se presenta de manera fija la etiqueta ``Total'' alineada al importe acumulado.

\begin{figure}[H]
    \centering
    \includesvg[width=0.4\textwidth]{images/cliente/CLI-15 — Cliente-Perfil.svg}
    \caption{Interfaz del perfil de usuario del cliente.}
    \label{fig:ui-cli-15}
\end{figure}

Esta interfaz corresponde a la pantalla de ``Perfil'', la cual contiene los datos requeridos para poder realizar pedidos en el sistema con información personal. En la sección superior se sitúa una barra de búsqueda con la etiqueta ``Busca en el menú'', ubicada a la derecha por el icono circular del carrito de compras. Inmediatamente abajo se incluye el control de navegación ``Regresar''.

El cuerpo central organiza los campos de datos personales del usuario mediante cajas de texto editables. Estas secciones corresponden a ``Nombre(s)'', ``Apellido(s)'', ``Correo(s)'', y el apartado de ``Contraseña'', la cual aparece encriptada con caracteres ocultos y dispone de un icono de ojo tachado en el extremo derecho para gestionar su visibilidad. En la base de la interfaz se mantiene fija la barra de navegación global que permite conmutar entre las pantallas de ``Menú'', ``Recolección'' y ``Perfil'', mostrando esta última pestaña seleccionada en color negro.

\begin{figure}[H]
    \centering
    \includesvg[width=0.4\textwidth]{images/cliente/CLI-16 — Cliente-Recoleccion-ResumenDelPedido-Anonimo.svg}
    \caption{Interfaz de recolección y resumen de pedido para usuario anónimo.}
    \label{fig:ui-cli-16}
\end{figure}

Esta interfaz corresponde a la pantalla de detalle de la orden en el apartado de recolección para un usuario que navega de forma anónima. La diferencia principal radica en la sección superior, donde se omite el botón de regreso y se introduce una barra de progreso horizontal en color azul acompañada por el texto estático ``Tu pedido se está preparando''. Inmediatamente abajo se expone el número de folio, seguido del botón de acción principal ``Escanear código'' configurado con un icono de código QR sobre un fondo negro, y el texto indicativo ``Realizado el'' que detalla la fecha y hora de la transacción.

El cuerpo de la pantalla mantiene la estructura estándar del flujo organizando la información en el bloque de ``Contenido'', el cual lista verticalmente los artículos adquiridos junto a sus respectivas fotografías, descripciones y cantidades en la base de cada tarjeta. En la sección posterior se detalla el ``Método de pago'' como ``Tarjeta de débito'' y el ``Horario seleccionado'' para la entrega. La interfaz concluye en el extremo inferior mostrando de manera fija la etiqueta ``Total'' alineada al importe acumulado.

\subsection{FLUJO DEL PERSONAL DE CAJA}